% Nouveau EGA V (cote 25): archive cote 25, site page 153 (PDF page 154).
% English translation of the checked French annotated master.
% AI-translated by GPT-6 Astra (OpenAI, ChatGPT Pro mode).
\scanpage{25/153}
\mspage{\hand{14}}

notions, apart from \var{\xout{\unsure{a}}}{}{typewriter strike-out}{ega25_25_153-v0001}\var{\inserted[typed]{the}}{the}{insertion}{ega25_25_153-v0002} dimensional condition = the condition of
$\underline{O}_X$-regularity, the character of x can be seen \var{\xout{\unsure{on}}}{}{typewriter strike-out}{ega25_25_153-v0003} in \unsure{fact} from the nature of
the point x in V(u), i.e. from the local ring $\underline{O}_{V(u),x}$ (possibly considered
as an algebra over k). When (in the case of a base field k)
x is rational over k, the notion relative to k coincides with the absolute
notion, as one sees \var{\inserted{by \unsure{proceeding|previously}}}{by \unsure{proceeding|previously}}{insertion}{ega25_25_153-v0004} as in the case where $\underline{E}$ is invertible. Finally,
when X is smooth over \var{\over{S}{Y}}{Y}{writing over}{ega25_25_153-v0005}, and u: $\underline{E}\to\underline{O}_X$ is again a homomorphism with
$\underline{E}$ loc free of finite type, one defines the corresponding notions relative
to S by considering fiber by fiber. Observe that if x is
nonsupersingular relative to \var{\over{S}{Y}}{Y}{writing over}{ega25_25_153-v0006}, then V(u) is flat over S \var{\xout{relativ}}{}{typewriter strike-out}{ega25_25_153-v0007} at
x (since u is $\underline{O}_X$-regular relative to \var{\over{S}{Y}}{Y}{writing over}{ega25_25_153-v0008} ...).

\var{\margin{$\Vert$}}{\placedmargin{$\Vert$}}{marginal note placed}{ega25_25_153-v0009} Of course (to be pointed out at least in a remark after the
formal definitions), one could just as well start from a section u
of a loc free Module $\underline{E}$; the preceding definitions apply by
interpreting u as a homomorphism $\overset{\hand{\scriptscriptstyle\vee}}{\underline{E}}\hand{\longrightarrow}\underline{O}_X$.

When X is smooth over \var{\over{S}{Y}}{Y}{writing over}{ega25_25_153-v0010}, u:$\underline{E}\hand{\longrightarrow}\underline{O}_X$ is given with $\underline{E}$ loc free
of rank m, then the set of \var{\xout{points}}{}{typewriter strike-out}{ega25_25_153-v0011} \var{\inserted[typed]{zeros of u}}{zeros of u}{insertion}{ega25_25_153-v0012} singular relative to S is the underlying
set of a well-defined closed subprescheme $V(u)^{\mathrm{sing}} = V^1$ of
$V(u)=V^o$, obtained as follows: consider $d^1u$, a section of $P^1_{X/\var{\over{S}{Y}}{Y}{writing over}{ega25_25_153-v0013}}(\overset{\hand{\scriptscriptstyle\vee}}{\underline{E}})$; its
restriction to $V^o$ can then be interpreted as a homomorphism
\[ \underline{E}\otimes\underline{O}_{V^o} \hand{\longrightarrow} \underline{\hand{\Omega}}^1_{X/\var{\over{S}{Y}}{Y}{writing over}{ega25_25_153-v0014}}\otimes\underline{O}_{V^o} \]
and $V(u)^{\mathrm{sing}}=V^1$ will denote the prescheme of zeros of \var{\xout{\unsure{the latter. Rem}}}{}{typewriter strike-out}{ega25_25_153-v0015}
\var{\xout{Considering similarly $d^2$, a section of $P^2_{X/S}(\overset{\hand{\scriptscriptstyle\vee}}{\underline{E}})$, then $d^2$ $V^1$, one}}{}{typewriter strike-out}{ega25_25_153-v0016}
\var{\xout{finds a canonical homomorphism}}{}{typewriter strike-out}{ega25_25_153-v0017} the corresponding homomorphism
\[ \hand{\bigwedge}^m\underline{E}\otimes\underline{O}_{V^o} \hand{\to} \underline{\hand{\Omega}}^m_{X/\var{\over{S}{Y}}{Y}{writing over}{ega25_25_153-v0018}}\otimes\underline{O}_{V^o} \]
i.e. of a well-defined section of $\underline{\hand{\Omega}}^m_{X/\var{\over{S}{Y}}{Y}{writing over}{ega25_25_153-v0019}}\otimes\mathrm{dét}(E)^{-1}\otimes\underline{O}_{V^o}$.

\mspage{\hand{153}}
\tnote{25-153-1}{Loc abbreviates locally. The doubt rendered on fact is the source's b/h in eb/eh fait. The inserted alternatives procédant/précédemment remain proceeding/previously, including the awkward second construction. The canceled fragments relativ and Rem remain incomplete. The source's several changes from S to Y do not replace the surviving S references in the surrounding prose. The formula's dét denotes determinant.}
